\section{Rendering --- target token architecture}
\label{sec:rendering-target}

This section is the \textbf{plan foundation} for finishing the migration.
Treat it as the working design; amend via the decision log.

\subsection{Goals}
\begin{enumerate}
  \item \textbf{API isolation:} only \pathref{Rendering/OpenGL/*} (and future
        backends) include GL/Vulkan/Metal headers for draw submission.
  \item \textbf{Stable front-end:} game modules describe \emph{what} to draw
        (meshes, textures, transforms, passes), not \emph{how} in GL terms.
  \item \textbf{Record then execute:} each frame builds a command stream
        (``tokens''), then the backend walks it once.
  \item \textbf{Resource handles:} CPU code holds opaque IDs; backend maps them
        to API objects.
  \item \textbf{Incremental:} keep the simulator playable after each step;
        prefer strangler-fig migration over a big bang rewrite.
\end{enumerate}

\subsection{Vocabulary}
\begin{description}[style=nextline]
  \item[Token / RenderCommand]
    A discrete instruction: bind X, set uniform Y, draw Z, begin pass, clear, \ldots
  \item[Command queue]
    Ordered buffer of tokens for the frame (or for a pass).
  \item[Backend]
    API-specific interpreter of tokens + owner of true GPU objects.
  \item[Enrollment]
    One-time (or rare) resource creation; returns a handle.
  \item[Front-end renderer]
    High-level helper that turns scene/draw requests into tokens
    (\symbolref{Renderer} today).
\end{description}

\subsection{Target frame pipeline}

\begin{center}
\begin{tikzpicture}[
  font=\small,
  box/.style={draw, rounded corners, align=center, inner sep=5pt, minimum width=2.4cm},
  arr/.style={-{Latex}, thick}
]
  \node[box, fill=orange!12] (mod) {Modules\\Update};
  \node[box, fill=orange!12, right=0.8cm of mod] (disp) {Dispatch /\\build draw list};
  \node[box, fill=purple!10, right=0.8cm of disp] (ren) {Renderer\\emit tokens};
  \node[box, fill=blue!10, right=0.8cm of ren] (q) {CommandQueue};
  \node[box, fill=green!12, below=0.9cm of q] (be) {IBackend\\Submit};
  \node[box, fill=gray!12, left=0.8cm of be] (gpu) {GPU API};

  \draw[arr] (mod) -- (disp);
  \draw[arr] (disp) -- (ren);
  \draw[arr] (ren) -- (q);
  \draw[arr] (q) -- (be);
  \draw[arr] (be) -- (gpu);
\end{tikzpicture}
\end{center}

\subsection{Recommended layering (target)}

\begin{lstlisting}[language={}]
Game / Rulesets / UI
    |  (no gl* includes)
    v
High-level draw requests / scene data
    |  (handles, CPU buffers, transforms)
    v
Renderer  ----enroll---->  IBackend::Create*
    |  (push tokens)
    v
CommandQueue<RenderCommand>
    |
    v
IBackend::SubmitCommandQueue
    |
    +-- GLBackend
    +-- (future) other backends
\end{lstlisting}

\subsection{Token payload design (proposal)}
\begin{openbox}
Current \symbolref{RenderCommand} is a single struct with many fields. Options:
\begin{enumerate}
  \item \textbf{Keep flat struct} for simplicity and POD friendliness (array
        queue, Tracy, no heap). Accept waste.
  \item \textbf{Tagged union / \texttt{std::variant}} for type safety and
        smaller conceptual surface.
  \item \textbf{Byte stream + opcodes} (more engine-y; harder to debug).
\end{enumerate}
Default recommendation for CSim scale: \textbf{(1) or a simple tagged union},
not a full GPU-driven byte compiler.
\end{openbox}

\subsection{What happens to Drawable?}
\begin{inferencebox}
Long-term, ``drawable'' should mean \emph{something that can append tokens}
(or register a renderable descriptor), not \emph{something that calls GL}.
CRTP can remain for CPU-side specialization, but GL handle fields should move
into the backend's handle tables.
\end{inferencebox}

Proposed end state for \symbolref{Canvas}:
\begin{itemize}
  \item Domain: \texttt{lifeCanvas}, dimensions, cell get/set (no GL).
  \item View/render: texture handle + mesh handle + material/pass ids;
        upload dirty regions via enrollment or an \texttt{UpdateTexture} token.
\end{itemize}

\subsection{Migration phases (strangler)}
\label{sec:migration-phases}

\begin{enumerate}[leftmargin=*]
  \item \textbf{Phase 0 --- Document \& freeze vocabulary} \\
        This document; agree command enum + handle meaning; stop adding new
        raw \texttt{gl*} call sites outside OpenGL backend / temporary
        drawables.

  \item \textbf{Phase 1 --- Make the queue real for one pass} \\
        Implement \symbolref{GLBackend::SubmitCommandQueue} for a minimal set:
        \symbolref{ClearColorBuffer}, \symbolref{SetViewport},
        \symbolref{SetShader}, \symbolref{SetTexture}, \symbolref{SetUniformMat4},
        \symbolref{DrawIndexed}.
        Unit-test or visually verify a fullscreen clear + one textured quad
        \emph{only} through tokens.

  \item \textbf{Phase 2 --- Route Scene through Renderer} \\
        \symbolref{Scene::Update} stops calling \texttt{glClear} / \symbolref{Draw}.
        Instead: build tokens (or ask each renderable to append tokens), then
        \symbolref{Renderer::RenderScene} $\rightarrow$ submit.
        Drawables may still be immediate behind a feature flag temporarily.

  \item \textbf{Phase 3 --- Migrate Canvas} \\
        Split domain vs GPU resources; canvas upload becomes handle-based;
        \symbolref{DrawImpl} becomes token emission.

  \item \textbf{Phase 4 --- Migrate remaining drawables} \\
        Splash text, GLString, command-line overlay, debug module, cursor, \ldots

  \item \textbf{Phase 5 --- Delete dead paths} \\
        Remove unused \symbolref{RenderQueue}-as-drawable-list if tokens won;
        collapse duplicate transform types; ensure Engine does not include GL.

  \item \textbf{Phase 6 --- Optional second backend} \\
        Only after OpenGL path is token-clean; e.g.\ mock backend for tests or
        a second real API.
\end{enumerate}

\begin{decisionbox}[title={Working order (proposed 2026-07-28)}]
Prefer \textbf{Phase 1 before large refactors of Game}. Prove the token
interpreter with a tiny known scene, then move Canvas. Do not invent Vulkan
until OpenGL is fully behind \symbolref{IBackend}.
\end{decisionbox}

\subsection{Acceptance criteria for ``migration done''}
\begin{itemize}
  \item No \texttt{gl*} / GLEW includes outside \pathref{Rendering/OpenGL/}
        (and maybe temporary test helpers).
  \item Frame submission path is enroll + token queue + submit.
  \item Cell sim still runs: edit, step/rules, camera, save/load hooks.
  \item \symbolref{Renderer::RenderScene} is non-empty and used.
  \item Decision log updated; dead stacks removed or quarantined under
        \pathref{archive/}.
\end{itemize}

\subsection{Risks}
\begin{itemize}
  \item Scope creep into full RHI (descriptor sets, multi-pass graph) before
        the sim needs them.
  \item Handle lifetime / deletion protocol undefined today.
  \item Performance paranoia too early (2048 command cap may be fine).
  \item Platform divergence (Linux/macOS files still older in places).
\end{itemize}
